\subsection{Firmware do nó sensor}
\label{subsec:firmware}

Os dois firmwares do nó, um para cada fluxo, foram escritos para o \gls{esp32cam} com o
\textit{core} Arduino-ESP32 3.0.7 e compartilham a mesma estrutura: configuração da câmera,
comunicação com o módulo de rádio, modo autônomo com \textit{deep sleep} e registro dos
tempos de cada etapa. Além do modo autônomo, ambos aceitam comandos pela porta serial,
usados nos testes e nas medições: disparo de uma leitura, ativação e desativação do modo
autônomo, \textit{deep sleep} avulso e repasse de comandos AT ao módulo de rádio.

\subsubsection{Ciclo de operação}
\label{subsubsec:ciclo_operacao}

No modo autônomo, o nó executa ciclos de leitura separados por períodos de
\textit{deep sleep}. Ao despertar pelo temporizador, o microcontrolador reinicia, configura
a câmera em tons de cinza na resolução de $160 \times 120$ pixels, captura um quadro e,
logo em seguida, desliga a câmera, pois ela não é necessária no restante do ciclo. A
partir desse ponto, cada fluxo segue um caminho, como resume a
Tabela~\ref{tab:ciclo_fluxos}.

\begin{table}[H]
    \centering
    \caption{Etapas de um ciclo de leitura em cada fluxo}
    \label{tab:ciclo_fluxos}
    \begin{tabularx}{\textwidth}{@{}lXX@{}}
        \toprule
        \textbf{Etapa} & \textbf{Fluxo em nuvem} & \textbf{Fluxo em borda} \\
        \midrule
        Despertar & \multicolumn{2}{l}{Reinício pelo temporizador e configuração da câmera} \\
        Captura & \multicolumn{2}{l}{Quadro de $160 \times 120$ pixels em cinza; câmera desligada em seguida} \\
        Processamento & Compressão JPEG e divisão em blocos de 175~bytes & Localização da janela
        e classificação dos dígitos com TensorFlow Lite Micro \\
        Envio & Um \textit{uplink} por bloco, com 7~s de intervalo (19 a 27 blocos) & Um único
        \textit{uplink} de 9~bytes \\
        Registro & \multicolumn{2}{l}{Linha de fim de ciclo na serial, com os tempos de cada etapa} \\
        Repouso & \multicolumn{2}{l}{\textit{Deep sleep} até o próximo ciclo} \\
        \bottomrule
    \end{tabularx}
    \legend{Fonte: Elaborado pelos autores (2026).}
\end{table}

Ao fim de cada ciclo, o firmware escreve na serial uma linha padronizada com o número do
ciclo, o resultado do envio e a duração de cada etapa (captura, processamento e envio).
Essa linha é lida pelos \textit{scripts} de medição, o que permite comparar os dois fluxos
com as mesmas ferramentas.

Antes de entrar em \textit{deep sleep}, o firmware desativa o driver da câmera, coloca o
sensor em modo de baixo consumo pelo pino de \textit{power-down} (GPIO~32), apaga o LED
de \textit{flash} (GPIO~4) e mantém o nível desses pinos durante o repouso. O contador de
ciclos e o estado do modo autônomo ficam na memória RTC, que é preservada durante o
\textit{deep sleep}. Apenas o ESP32 é colocado em repouso: o módulo Radioenge continua
alimentado e não recebe comando para entrar em modo suspenso, o que mantém nele a sessão
\gls{lorawan} e o contador de quadros, mas também mantém o seu consumo durante o repouso,
como discutido nos resultados. Nas medições, o intervalo de \textit{deep sleep} foi de
30~s nos dois fluxos.

\subsubsection{Execução dos modelos no microcontrolador}
\label{subsubsec:tflm}

No fluxo em borda, os modelos são executados com o TensorFlow Lite Micro, na versão
\texttt{esp-tflite-micro} 1.3.1 com a biblioteca de otimização \texttt{esp-nn} 1.1.0, ambas
incluídas no \textit{core} 3.0.7. Os dois modelos \texttt{.tflite} ficam na memória
\textit{flash} como vetores de bytes, e cada um tem o seu próprio interpretador, criado com
um registro de operadores que contém apenas os operadores usados pelo modelo: sete na
localização da janela e oito no classificador. O firmware compilado ocupa 563~KB, cerca de
17\% da memória de programa.

O TensorFlow Lite Micro não aloca memória dinamicamente: cada interpretador recebe uma
área de trabalho (\textit{arena}) de tamanho fixo, onde ficam os tensores intermediários.
As duas \textit{arenas} foram alocadas na PSRAM de 4~MB da placa. A
Tabela~\ref{tab:memoria_tflm} mostra a memória reservada e a efetivamente usada, medida na
placa durante 61 capturas.

\begin{table}[H]
    \centering
    \caption{Uso da PSRAM pelo fluxo em borda}
    \label{tab:memoria_tflm}
    \begin{tabular}{@{}lrr@{}}
        \toprule
        \textbf{Item} & \textbf{Reservado (bytes)} & \textbf{Usado (bytes)} \\
        \midrule
        Buffer de quadro da câmera & --- & $\approx$19.216 \\
        \textit{Arena} da localização da janela & 2.097.152 & 1.623.852 \\
        \textit{Arena} do classificador de dígitos & 786.432 & 28.556 \\
        PSRAM livre após as alocações & \multicolumn{2}{r}{1.289.292} \\
        \bottomrule
    \end{tabular}
    \legend{Fonte: Elaborado pelos autores (2026).}
\end{table}

O dimensionamento das \textit{arenas} exigiu ajustes durante os testes em placa. A
primeira estimativa, baseada no maior tensor intermediário do modelo de localização,
reservou 1~MB, e a alocação falhou: o interpretador mantém vários tensores intermediários
ativos ao mesmo tempo, e o modelo precisou de cerca de 1,6~MB. A \textit{arena} foi então
ampliada para 2~MB. A \textit{arena} do classificador, de 768~KB, mostrou-se
superdimensionada, pois o modelo usa menos de 29~KB. Para evitar novas tentativas na
placa, foi escrito um \textit{script} que reproduz o planejador de memória do
TensorFlow Lite Micro a partir do arquivo \texttt{.tflite}; para o modelo de localização,
ele estimou 1.622.016~bytes, contra 1.623.852~bytes medidos. Esse mesmo \textit{script}
confirmou que o modelo de localização do pipeline A, que pediu 15,3~MB, não caberia na
placa.

A leitura no firmware segue os passos descritos na
Seção~\ref{subsubsec:leitura_roi}. O quadro capturado é copiado para um buffer próprio e a
câmera é desligada antes da inferência. Para o modelo de localização, o quadro é
redimensionado por vizinho mais próximo para $192 \times 192$ pixels diretamente no
tensor de entrada, já convertido para INT8 pela subtração de 128 de cada pixel. Na saída,
os valores INT8 maiores que zero, que correspondem a probabilidade superior a 0,5, formam
a máscara, e a caixa é o menor retângulo que a contém. Cada uma das $N$ fatias da caixa é
redimensionada para $28 \times 28$ pixels e classificada, e o dígito de maior pontuação é
escrito na leitura. Se a máscara vier vazia, a leitura é enviada com traços no lugar dos
dígitos.

O tempo de inferência medido na placa foi muito maior do que o esperado para modelos tão
pequenos. Em 61 capturas, a localização da janela levou mediana de 15.564~ms, com variação
de apenas 4~ms entre capturas, e a classificação dos seis dígitos levou 1.995~ms, para um
total de 17.838~ms por leitura. A causa não foi investigada; uma hipótese é que parte dos
operadores do modelo de localização, como a convolução transposta, não tenha versão
otimizada na \texttt{esp-nn} e seja executada pelas implementações de referência, em C
genérico. Durante essas capturas, a temperatura interna do chip permaneceu entre 72 e
74~°C.

Como discutido na Seção~\ref{subsubsec:leitura_roi}, a solução adotada dispensa o modelo de
localização e usa uma ROI calibrada na instalação. O firmware usado nas medições, porém,
ainda executa a localização a cada leitura; a versão com ROI fixa não foi implementada.
Sem a localização, a leitura no nó ficaria restrita à classificação dos dígitos, de cerca
de 2~s, o que é uma estimativa, e não uma medição.

\subsubsection{Firmware de depuração e calibração}
\label{subsubsec:firmware_debug}

Um terceiro firmware, sem rádio, foi usado para posicionar a câmera diante do medidor e
para medir o desempenho dos modelos na placa. A cada 5~s, ele captura um quadro, executa a
leitura completa, desenha sobre a imagem a caixa encontrada e as divisões das fatias e
envia ao computador a imagem anotada, a leitura e métricas de memória, tempo e
temperatura. Um \textit{script} no computador grava as imagens e as métricas em arquivo.
Com ele, a câmera foi ajustada à mão até que a janela de dígitos ocupasse o quadro e a
caixa prevista coincidisse com ela; a partir desse ponto, a leitura do medidor de teste se
estabilizou. Os dados de memória e de tempo de inferência apresentados nesta seção vêm
desse firmware.

\subsection{Servidor da aplicação}
\label{subsec:servidor}

No servidor, a recepção é feita pelo \textit{script} \texttt{reconstrutor.py}, descrito na
Seção~\ref{subsubsec:recepcao}: no fluxo em borda, a leitura chega pronta na porta~2; no
fluxo em nuvem, a imagem remontada é gravada em disco. O \gls{ocr} do fluxo em nuvem é
executado sobre essas imagens pelo \textit{script} \texttt{pipeline\_nuvem.py}.

Para que a única diferença entre os fluxos seja o local do processamento, o servidor
executa os mesmos arquivos \texttt{.tflite} INT8 do firmware, com o interpretador LiteRT
2.3.0, e reproduz a mesma lógica: o mesmo redimensionamento por vizinho mais próximo, com
o mesmo cálculo de índices, a mesma conversão para INT8, o mesmo limiar na máscara e o
mesmo fatiamento. O acelerador XNNPACK do LiteRT foi desativado, para que os operadores
sejam executados pelas implementações de referência, como no microcontrolador. A única
diferença no dado de entrada é a compressão JPEG com qualidade 90, que o fluxo em borda não
sofre.

\subsubsection{Tempo de inferência no servidor}
\label{subsubsec:tempo_servidor}

Na implementação atual, o \gls{ocr} do servidor não é disparado automaticamente quando a
imagem fica completa: ele é executado depois, sobre as imagens gravadas. Por isso, a
latência medida no fluxo em nuvem termina no momento em que a imagem fica disponível no
servidor, e o tempo de inferência foi medido à parte. Em um computador com processador
AMD Ryzen 5 5600GT, as 21 imagens completas recebidas pelo enlace foram lidas cinco vezes
cada. A Tabela~\ref{tab:tempo_inferencia} compara os tempos com os do microcontrolador.

\begin{table}[H]
    \centering
    \caption{Tempo de inferência por leitura no nó e no servidor (mediana)}
    \label{tab:tempo_inferencia}
    \begin{tabular}{@{}lrr@{}}
        \toprule
        \textbf{Etapa} & \textbf{\gls{esp32cam} (ms)} & \textbf{Servidor (ms)} \\
        \midrule
        Localização da janela & 15.564 & 4,30 \\
        Leitura completa (localização e seis dígitos) & 17.838 & 6,25 \\
        Leitura com ROI fixa (só os seis dígitos) & $\approx$1.995 & 1,91 \\
        \bottomrule
    \end{tabular}
    \legend{Fonte: Elaborado pelos autores (2026). No servidor, os tempos incluem a
    decodificação do JPEG; no \gls{esp32cam}, a leitura com ROI fixa corresponde ao tempo
    medido da classificação dos dígitos.}
\end{table}

O servidor executa a leitura completa cerca de 2.800 vezes mais rápido que o
microcontrolador. Diante da latência de envio da imagem, de mais de dois minutos, o tempo
de inferência no servidor é desprezível, e a omissão dessa etapa na latência medida não
altera a comparação entre os fluxos.

\subsubsection{Recursos disponíveis apenas no servidor}
\label{subsubsec:recursos_servidor}

O baixo custo da inferência no servidor permite etapas que seriam caras no
microcontrolador. Duas foram implementadas e avaliadas:

\begin{itemize}
    \item \textbf{Correção da inclinação}: a fileira de roletes costuma aparecer levemente
    inclinada no quadro. O servidor detecta as bordas da imagem (Canny), extrai os segmentos
    de reta pela transformada de Hough, toma a mediana ponderada pelo comprimento dos
    ângulos dos segmentos quase horizontais (até $\pm 20^\circ$) e gira o quadro para
    endireitá-lo antes da leitura.
    \item \textbf{Votação entre caixas deslocadas}: para tolerar pequenos erros na posição
    da ROI, o classificador é aplicado a cinco versões da caixa, a original e outras quatro
    deslocadas de 2~pixels na horizontal e de 1~pixel na vertical, e as pontuações de cada
    fatia são somadas antes da escolha do dígito. No microcontrolador, isso multiplicaria
    por cinco os cerca de 2~s da classificação.
\end{itemize}

A ROI é informada ao servidor como um retângulo em coordenadas do quadro de
$160 \times 120$ pixels e é calibrada uma vez por instalação. Para isso, um
\textit{script} auxiliar desenha uma grade de coordenadas sobre o quadro ampliado, a
partir da qual a janela é lida à mão. Ainda não existe uma ferramenta que faça essa
calibração de forma automática ou guiada, o que é uma limitação da solução atual. O efeito
da correção da inclinação, da votação e da precisão da ROI sobre a acurácia é apresentado
nos resultados.
